\section{What is saved, and what preparation still costs}
\label{sec:costs}

A small deployed state, a cheap initialization, and a cheap query are
different claims. Legendre starts from the supplied dense initialization
and creates its moments directly. It saves changing memory but still
multiplies by the fixed dense matrices, so it does not remove their
quadratic cost. The selected models discard those matrices after setup.
Their nonlinear forward and backward passes then use only the retained
small matrices, metrics and training data. No width-$n$ oracle is called
during their subsequent evolution.

\paragraph{The common selected runtime.}
Both source builders lead to the same execution pattern: a small forward
pass, metric-adjoint backward pass, and a training-feature Gram solve for
the readout correction. The solve is part of training, not an uncounted
oracle. Refreshing the corrected readout once per state permits later
queries to use a plain small forward pass. Passive labels are absent from
all these operations. For Logarithmic, the certified passive queries
were declared during setup; the existence of an evaluable network on other
inputs is not a guarantee of accuracy there.

\paragraph{Information available at initialization.}
Analyticity makes every finite coefficient list recoverable, to a prescribed
tolerance, from finitely many derivatives at time zero. Those derivatives
are obtained by differentiating the dense training vector field at its
initial weights, using the training labels. Passive inputs only specify
which response curves are represented. The continuation procedure in
\Cref{sec:panel-proof} reconstructs the later polynomial coefficients
without access to later trained weights. This is a data-dependent
compression of a specified learning problem, not a data-oblivious map.

That statement is about available information, not computational economy.
Analytic continuation from one point can require a very large derivative
order and high precision. A practical dense rollout can supply source
coefficients more conveniently, but its work and full-horizon information
must then be reported as setup. Such an implementation is not evidence
for a cheap initialization-only algorithm.

\paragraph{Which efficient setup bounds are already available.}
Harmonic has two separately certified implementations. At fixed problem
parameters, explicit dense continuation costs near-quadratic arithmetic
work; implicit Gaussian continuation costs near-linear work by revealing
only the random matrix actions actually needed. The latter samples a fresh
coupled dense reference: it cannot read a supplied dense matrix in
subquadratic time. Their temporary arrays are discarded before deployment.

The new Logarithmic theorem does not inherit either setup-time bound merely
because its runtime is the same. Its enlarged source domain and finite-jet
compiler prove the cubic-logarithmic retained state, not a corresponding
fast preprocessing algorithm. \Cref{app:execution-costs} preserves the
detailed existing execution ledgers, and
\hyperref[int:harmonic-efficient-setup-proofs]{the setup appendix}
proves the two Harmonic bounds. Real-coordinate arithmetic, evaluator
cost, numerical precision and the number of integration steps remain
separate resources throughout.
